% Stile MetiX per il PDF di Sphinx: testo in Pagella, titoli e apparato in sans (TeX Gyre Heros),
% nero per titoli e testo, arancione solo per gli accenti grafici (mai per il testo).
% Caricato in fondo al preambolo di Sphinx (latex_elements['preamble'] in conf.py).
% Palette MetiX: nero per titoli e testo, arancione solo per gli accenti grafici (filetti, barre,
% pallini, logo), mai per il testo: sul bianco l'arancione si legge male.
\definecolor{metixarancio}{HTML}{E8590C}
\definecolor{metixnero}{HTML}{18181B}
\definecolor{metixinchiostro}{HTML}{27272A}
\definecolor{metixgrigio}{HTML}{71717A}
\definecolor{carta}{HTML}{F4F4F5}
% etichette di origine e di stato
\definecolor{nostro}{HTML}{15803D}
\definecolor{letteratura}{HTML}{2563EB}
\definecolor{misto}{HTML}{0D9488}
\definecolor{sospeso}{HTML}{B91C1C}
\definecolor{sperim}{HTML}{7C3AED}
% nomi usati dalle figure TikZ: tratti strutturali e accenti
\colorlet{metixpetrolio}{metixnero}  % nome storico
\colorlet{metixblu}{metixnero}
\colorlet{metixazzurro}{metixarancio}

\usepackage{microtype}
\usepackage{csquotes}
\usepackage{mathtools}
\usepackage{caption}
\usepackage{enumitem}
\usepackage{tikz}
\usetikzlibrary{arrows.meta,positioning,shapes.geometric,calc,fit,backgrounds,decorations.pathreplacing}
\usepackage[skins,breakable]{tcolorbox}
\usepackage[chapter]{algorithm}
\usepackage{algorithmic}
\usepackage{siunitx}
\sisetup{output-decimal-marker={,},group-separator={\,},group-minimum-digits=4,per-mode=symbol}

\newfontfamily\metixsans{texgyreheros}[Extension=.otf,UprightFont=*-regular,BoldFont=*-bold,
  ItalicFont=*-italic,BoldItalicFont=*-bolditalic,Scale=0.92]
\AtBeginDocument{\color{metixinchiostro}}

% ---------------------------------------------------------------- titoli
\titleformat{\chapter}[display]
  {\normalfont\metixsans}
  {\color{metixgrigio}\small\bfseries\addfontfeatures{LetterSpace=15}%
   \MakeUppercase{\chaptertitlename}\enspace\thechapter}
  {8pt}
  {\color{metixnero}\fontsize{30}{34}\selectfont\bfseries}
  [{\color{metixarancio}\rule{2.4cm}{2.5pt}}]
\titlespacing*{\chapter}{0pt}{-6pt}{28pt}
\titleformat{\section}{\normalfont\metixsans\Large\bfseries\color{metixnero}}
  {\color{metixgrigio}\thesection}{0.7em}{}
\titleformat{\subsection}{\normalfont\metixsans\large\bfseries\color{metixnero}}
  {\color{metixgrigio}\thesubsection}{0.7em}{}
\titlespacing*{\section}{0pt}{3.2ex plus 1ex minus .2ex}{1.4ex plus .2ex}
\newcommand{\metixpar}[1]{{\metixsans\bfseries\color{metixnero}#1}\enspace}

% ---------------------------------------------------------------- testatine
\fancypagestyle{normal}{%
  \fancyhf{}%
  \fancyhead[L]{\raisebox{-0.2ex}{\includegraphics[height=2.2mm]{marchio-inchiostro.pdf}}\enspace
    \metixsans\footnotesize\color{metixgrigio}manuale tecnico}%
  \fancyhead[R]{\metixsans\footnotesize\color{metixgrigio}\nouppercase{\leftmark}}%
  \fancyfoot[C]{\metixsans\footnotesize\color{metixgrigio}\thepage}%
  \renewcommand{\headrulewidth}{0pt}%
  \renewcommand{\footrulewidth}{0pt}}
\fancypagestyle{plain}{%
  \fancyhf{}\fancyfoot[C]{\metixsans\footnotesize\color{metixgrigio}\thepage}%
  \renewcommand{\headrulewidth}{0pt}\renewcommand{\footrulewidth}{0pt}}
\renewcommand{\chaptermark}[1]{\markboth{\ifnum\value{chapter}>0\thechapter\enspace\fi#1}{}}

% ---------------------------------------------------------------- testo
\setlist{itemsep=2pt,topsep=4pt,leftmargin=1.4em}
\setlist[itemize]{label={\color{metixarancio}\small\textbullet}}
\setlength{\emergencystretch}{2.5em}
\DeclareCaptionFont{metixnero}{\color{metixnero}}
\captionsetup{font={small,sf},labelfont={bf,metixnero},format=hang,
  justification=raggedright,singlelinecheck=true}
\renewcommand{\arraystretch}{1.18}

% etichette di origine e di stato: pillole sans
\newcommand{\etichetta}[2]{\tikz[baseline=(t.base)]\node[fill=#1,text=white,rounded corners=3pt,
  inner xsep=3.5pt,inner ysep=1.6pt,font=\metixsans\scriptsize\bfseries](t){#2};}
\newcommand{\Nostro}{\etichetta{nostro}{NOSTRO}}
\newcommand{\Lett}{\etichetta{letteratura}{LETTERATURA}}
\newcommand{\Misto}{\etichetta{misto}{LETT.\,+\,NOSTRO}}
\newcommand{\Sospeso}{\etichetta{sospeso}{SOSPESO}}
\newcommand{\Sperim}{\etichetta{sperim}{SPERIMENTALE}}
\newcommand{\Opz}{\etichetta{metixgrigio}{OPZIONE}}
% stato dell'implementazione: dove vive il codice
\newcommand{\etichettachiara}[2]{\tikz[baseline=(t.base)]\node[draw=#1,line width=0.6pt,text=#1,
  rounded corners=3pt,inner xsep=3.2pt,inner ysep=1.3pt,font=\metixsans\scriptsize\bfseries](t){#2};}
\newcommand{\StatoRust}{\etichetta{metixnero}{RUST}}
\newcommand{\StatoPython}{\etichettachiara{metixnero}{SOLO PYTHON}}
\newcommand{\StatoContorno}{\etichettachiara{metixgrigio}{CONTORNO}}

% riquadri: fondo tenue e barra laterale nel colore del riquadro, titolo sans nero spaziato
\tcbset{metixriquadro/.style={enhanced,breakable,frame hidden,boxrule=0pt,arc=0pt,
  left=10pt,right=8pt,top=5pt,bottom=5pt,
  fonttitle=\metixsans\scriptsize\bfseries\addfontfeatures{LetterSpace=10},
  coltitle=metixnero,colbacktitle=#1!6,colback=#1!6,
  attach title to upper={\par\smallskip},
  borderline west={3pt}{0pt}{#1}}}
\newtcolorbox{origine}[1][]{metixriquadro=metixarancio,title={ORIGINE DEL METODO},#1}
\newtcolorbox{nota}[1][]{metixriquadro=metixgrigio,#1}
\newtcolorbox{attenzione}[1][]{metixriquadro=sospeso,title={ATTENZIONE},#1}
\newtcolorbox{misura}[1][]{metixriquadro=nostro,title={MISURA},#1}
\newtcolorbox{implementazione}[1][]{metixriquadro=metixnero,colback=metixnero!3,colbacktitle=metixnero!3,
  title={IMPLEMENTAZIONE},#1}

% codice in linea: più piccolo, può andare a capo dopo _
\newcommand{\spezza}{\discretionary{}{}{}}
\DeclareRobustCommand{\ttbreak}{\def\_{\textunderscore\spezza}}
\let\metixorigcode\sphinxcode
\renewcommand{\sphinxcode}[1]{{\small\color{metixnero}\ttbreak\metixorigcode{#1}}}
\newcommand{\file}[1]{{\small\ttfamily\ttbreak#1}}
\newcommand{\vect}[1]{\symbf{#1}}
\newcommand{\R}{\mathbb{R}}

% algoritmi in italiano
\renewcommand{\algorithmicrequire}{\textbf{Dati:}}
\renewcommand{\algorithmicensure}{\textbf{Risultato:}}
\renewcommand{\algorithmicfor}{\textbf{per}}
\renewcommand{\algorithmicdo}{\textbf{:}}
\renewcommand{\algorithmicendfor}{\textbf{fine}}
\renewcommand{\algorithmicif}{\textbf{se}}
\renewcommand{\algorithmicthen}{\textbf{allora}}
\renewcommand{\algorithmicendif}{\textbf{fine}}
\floatname{algorithm}{Algoritmo}

% tabelle: intestazioni sans, tabelle lunghe in corpo minore
\renewcommand{\sphinxstyletheadfamily}{\metixsans\bfseries\color{metixnero}}
\AtBeginEnvironment{longtable}{\small}
